% C-counts.tex: числа и то, как каждое измерено.
\section{Величины на 1.0.0-кв.7 и как они измерены}
\label{app:counts}

Каждая величина в настоящем документе снята с дерева на версии \code{1.0.0-кв.7} (коммит
\code{4e74376}, 20~сентября 2026 года) способом, указанным рядом с ней, а не переписана из более
раннего документа. Там, где собственное отмеченное число репозитория измерено на более ранней
версии, эта версия названа. Там, где число внешнее, строка говорит, чьё оно.

Версия~3 настоящего документа была написана на \code{1.0.0-кв.1} и перемерена здесь.
Четырнадцать из двадцати шести величин ниже не изменились; сдвинувшиеся строки отмечены в тексте
там, где это важно. Изменилось не одно число, а один способ: число маршрутов прежнее, а команда,
которая его давала, уже нет.

{\footnotesize
\setlength{\tabcolsep}{4pt}\renewcommand{\arraystretch}{1.12}
\rowcolors{2}{tabstripe}{white}
\begin{xltabular}{\textwidth}{>{\RaggedRight\arraybackslash\hspace{0pt}}p{2.1in} >{\raggedleft\arraybackslash}p{0.78in} >{\RaggedRight\arraybackslash\hspace{0pt}}X}
\caption{Величины на \code{1.0.0-кв.7} (коммит \code{4e74376}) и способ их получения. Дерево меняется и после измерения, так что величина, снятая позже, может отличаться; воспроизводимой каждую делает именно коммит.}\label{tab:counts}\\
\toprule
\rowcolor{tabheader}\thead{Величина} & \thead{Значение} & \thead{Способ} \\
\midrule
\endfirsthead
\toprule
\rowcolor{tabheader}\thead{Величина} & \thead{Значение} & \thead{Способ} \\
\midrule
\endhead
\midrule\multicolumn{3}{r}{\itshape продолжение на следующей странице}\\
\endfoot
\bottomrule
\endlastfoot
Канонические таблицы в \code{01\_схема.скл} & 45 & Файл объявляет 50 инструкций \code{СОЗДАТЬ ТАБЛИЦУ}, пять из которых суть дочерние разделы по умолчанию для разбитых таблиц событий. Мигрированное развёртывание держит 52 таблицы. \\
Триггеры в \code{06\_триггеры.скл} & 42 & Число строк \code{СОЗДАТЬ ТРИГГЕР} \\
Триггеры в загруженной схеме & 46 & 42 из файла и четыре, объявленные рядом с тем, что они охраняют: один в файле реестра миграций, один в \code{01\_схема.скл}, два в \code{12\_в7\_ограничения.скл}. Прочитано из системного каталога триггеров загруженной и мигрированной базы, без внутренних триггеров и без 25 копий, которые ПостгреСКЛ создаёт на разделах разбитых таблиц \\
Таблицы под триггером добавления & 31 & Триггеры \code{ДО ОБНОВЛЕНИЯ ИЛИ УДАЛЕНИЯ} из \code{06\_триггеры.скл} и миграций, по таблицам; проектная запись называет среди них четырнадцать случаев аудита как самой записи \\
Инвариантные проверки & 313 & Число строк, начинающихся с \code{проверка\_} в определении функции, в \code{полярис\_проверки/\allowbreak{}проверки.пи}, и оно равно зарегистрированному перечню; у каждой есть тест обнаружения с положительным контролем в \code{тест\_проверки.пи} \\
Мутационные учения & 12 & Файлы вида \code{сценарии/\allowbreak{}полярис-*мутационное-учение*}: над проверками, ограничениями, триггерами, процедурами, свидетелями нулевого разглашения, контрактом соответствия, двумя КСР, верификатором ОИД4ВП, приложением, покрытием, конституцией и перемещениями между модулями \\
Маршруты приложения & 124 & Число объявлений маршрутов во всех файлах \code{полярис\_веб/*.пи}. Число не изменилось с \code{1.0.0-кв.1}, а СПОСОБ изменился. 18~сентября 2026 года \code{прил.пи} был разложен на десять модулей маршрутов, и прежний способ, читавший один этот файл, теперь даёт 13 \\
Задачи НИ в \code{ни.ямл} & 20 & Ключи задач рабочего процесса; ещё два рабочих процесса выполняются ежемесячно и еженедельно, а один, запускаемый вручную, публикует пакеты \\
Подписанные форматы передачи в спецификации & 27 & Различные строки \code{полярис-*/1} в \code{документы/\allowbreak{}справка/\allowbreak{}СПЕЦИФИКАЦИЯ-ПЕРЕДАЧИ.мд} \\
Случаи соответствия & 118 & Случаи файла \code{соответствие/случаи.джсон} по двадцати пяти типам предметов; замороженный набор версии 1 держит 44 \\
Поля вердикта, которых не стесняет ни один опубликованный случай & 44 & Объявленный список выживших в мутационном учении над контрактом соответствия, который учение держит равным собственному измерению в обе стороны: из 89 полей решения у 27 верификаторов, до которых доходят случаи, 44 переживают принудительное разрешение, и ни одно из них нельзя стеснить, написав случай \\
Формальные модели, проверяемые в НИ & 4 & \code{мета/тла/*.тла}; у трёх есть парная конфигурация, на которой они обязаны упасть \\
Сценарии учений & 80 & Файлы вида \code{сценарии/\allowbreak{}полярис-*учение.*}, среди них двенадцать мутационных \\
Миграции & 45 & Файлы вида \code{полярис\_скл/\allowbreak{}миграции/\allowbreak{}*.вверх.скл}, у каждой есть откат \\
Руководства и книги оператора & 21 & Файлы разметки в \code{документы/оператор/}, включая указатель \\
Помеченные версии & 3 & Метки репозитория. На \code{1.0.0-кв.1} их было 295, по одной на каждую отправку от \code{в9.30} (16~мая 2026 года) до \code{в9.466} (15~сентября 2026 года). Метка отмечает ВЫПУСК, а не отправку, поэтому те были сняты; остались \code{в1.0.0-кв.1}, \code{-кв.2} и \code{-кв.3}, три опубликованных кандидата \\
Определённые функции тестов & 2\,189 & Число определений функций \code{тест\_} в наборах продукта, командной строки, слоя проверок, сценариев, карты, моделирования, КСР и пакетов; это число определений, а не прогонов \\
Проходящие тесты продукта, по отметке в ЧИТАЙМЕНЯ & 1\,030 & Измерено на \code{1.0.0-кв.7} на эталонной машине, по отдельному прогону каждого набора; три пропускаются без настоящего МЛ-ДСА \\
Тесты криптографических свидетелей, по отметке & 107 из 112 & Измерено на \code{1.0.0-кв.7}, без изменений с \code{1.0.0-кв.1}; пяти нужен модуль ПККС\#11 или настоящий ключ СУК, и НИ прогоняет все, кроме того, что требует СУК \\
Тесты пакета верификатора ОИД4ВП & 244 & Прогон в \code{пакеты/полярис-оид4вп}; три из них суть положительный контроль, без которого отказы были бы пусты \\
Пакеты в публичных реестрах & 4 & \code{полярис-проверка}, \code{полярис-оид4вп}, \code{полярис-сдк-питон} в реестре ПайПИ; \code{полярис-сдк-тс} в реестре НПМ; каждый проверен установкой из живого реестра. Когда писалась версия~3, все четыре отдавали \code{1.0.0-кв.1}. \code{1.0.0-кв.3}, опубликованный 18~сентября 2026 года, есть предварительный выпуск: установщик берёт его, лишь когда об этом попросили (\code{пип установить --пред}, метка НПМ \code{следующий}), а обычная установка по-прежнему даёт \code{0.1.0}, последний стабильный выпуск. Считано с живых реестров 23~сентября 2026 года \\
Размещённые модули проверки соответствия & 11 & Служба Фонда ОпенАйДи, 15~сентября 2026 года: 7 \code{ПРОЙДЕНО}, 4 \code{НА РАССМОТРЕНИИ}, ноль отказов, ноль предупреждений; собственные подписанные выгрузки службы \\
Внешние векторы проверки МЛ-ДСА-65 & 58 из 58 & Проект Вайчпруф, закреплённый на коммите, под обоими свидетелями при каждой отправке \\
Названные сторонние реализации, выполнившие предъявление & 1 & Уолт.ид Воллет ППИ в2 по закреплённому дайджесту образа, 15~сентября 2026 года; один путь, один формат \\
Исправления, вызванные внешней находкой & 2 & Строка сводной таблицы: сертификат, который отверг кошелёк (15~сентября 2026 года), и отметка версии, которую сторонний рецензент не смог прочесть однозначно (17~сентября 2026 года) \\
Независимые ревью безопасности, пилоты, операторы, не являющиеся автором & 0 & Строки сводной таблицы, пустые \\
\end{xltabular}
}

Числа производительности в настоящем документе суть числа таблицы~\ref{tab:perf}, каждое с
версией, которая его измерила. Базовая линия приложения не перезапускалась с \code{в9.191} и
приводится с этой отметкой. Стенд производительности перезапущен 23~сентября 2026 года на дереве
\code{1.0.0-кв.7}, с доказывающей программой, пересобранной из него, и перезапуск обнаружил, что его
измерение доказательства с \code{в9.354} не измеряло ничего: он по-прежнему посылал доказывающей
программе вход, которого та больше не принимала, печатал \emph{здесь не измерено} и завершался
чисто, так что его шаг НИ оставался зелёным. Строки подписей сдвинулись на несколько процентов;
строка доказательства упала с 35~мс до 25~мс, и сравнение не равноценно, потому что доказываемое
утверждение изменилось на \code{в9.354}. Половина, которая не измерила, теперь роняет шаг. Эталонная машина во
всём репозитории есть ноутбук на процессоре Эппл с восемью ядрами и шестнадцатью гигабайтами
памяти.
